\section{Spec 5 --- Information extraction}

We extract six entity types (water body, location, organism, quantity, date/time, weather) and twelve water conditions
with three systems: spaCy pipelines for 18 languages~\cite{honnibal2020spacy}, a multilingual XLM-RoBERTa NER
model~\cite{conneau2020xlmr}, and our rules (water-term lists in nine languages and regular expressions). The NER
models know places, dates and numbers but not organisms, water bodies or weather, which only the rules cover. The
same NER output was used to mask personal names: 8.8\% of notes named a person, mostly survey credits.

\section{Spec 6 --- Gold set}

We labelled 300 records from one stratified sample, plus 20 calibration records labelled by
all three of us. Each gold record was labelled by one student and one LLM
(Qwen2.5-VL-7B~\cite{bai2025qwen25vl}) and a third person decided the disagreements. Table~\ref{tab:gold} summarises the
labels.

\begin{table}[htbp]
\centering\small
\caption{Gold set (300 records).}
\label{tab:gold}
\resizebox{\columnwidth}{!}{%
\begin{tabular}{ll}
\toprule
Label & Distribution \\
\midrule
Relevance & about water 150, not 134, uncertain 16 \\
Conditions & none stated 128; low or dry 16, high water 7, clear 6, turbid 4 \\
Entities (144 spans) & water body 54, quantity 28, weather 18, location 16, organism 15, date 13 \\
Image--text & consistent 94, cannot judge 11, inconsistent 5, partly 3 \\
\bottomrule
\end{tabular}}
\end{table}

\section{Spec 7 --- Agreement and extractor evaluation}

The agreement is measured with Cohen's $\kappa$~\cite{cohen1960kappa} and span F1 (Table~\ref{tab:agreement}). 
The agreement between us students was high, but the agreement with the LLM is low - 68\% of notes got marked as ``not about water'', 
while we students only labelled 40\% this way.

\begin{table}[htbp]
\centering\small
\caption{Agreement: relevance and condition $\kappa$, span F1 (exact / partial).}
\label{tab:agreement}
\resizebox{\columnwidth}{!}{%
\begin{tabular}{lrrr}
\toprule
Pair & Relevance $\kappa$ & Conditions $\kappa$ & Span F1 \\
\midrule
Students, calibration (20) & 0.80 & 1.00 & 1.00 / 1.00 \\
Student vs LLM, gold (300) & 0.25 & 0.29 & 0.20 / 0.29 \\
\bottomrule
\end{tabular}}
\end{table}

The extractors are scored on the gold entity spans (Table~\ref{tab:extractors}). Exact scores are low because span
boundaries differ from the guideline. With partial matching the rules are best.
For water conditions the rules reach F1 0.31 (precision 0.50, recall 0.22). Only 51\% of the
gold notes are actually about water. Translating non-English notes to English helps condition extraction but
hurts entity extraction, because translation garbles names and abbreviations (Table~\ref{tab:mt}).

\begin{table}[htbp]
\centering\small
\caption{Extractors on the gold set (150 spans, micro F1).}
\label{tab:extractors}
\begin{tabular}{lrrr}
\toprule
System & F1 exact & F1 partial & Predicted \\
\midrule
spaCy & 0.20 & 0.34 & 96 \\
XLM-R NER & 0.20 & 0.28 & 65 \\
Rules & 0.12 & \textbf{0.44} & 210 \\
Union & \textbf{0.22} & 0.40 & 332 \\
\bottomrule
\end{tabular}
\end{table}

\begin{table}[htbp]
\centering\small
\caption{Extraction F1 on English originals vs non-English notes, original and machine-translated.}
\label{tab:mt}
\resizebox{\columnwidth}{!}{%
\begin{tabular}{lrrr}
\toprule
Text & spaCy entities & Rules entities & Rules conditions \\
\midrule
English originals & 0.42 & 0.53 & 0.39 \\
Non-English, original & 0.28 & 0.37 & 0.22 \\
Non-English, translated & 0.18 & 0.32 & 0.37 \\
\bottomrule
\end{tabular}}
\end{table}
